\todo{只提取与本研究最相关的大约30-40篇Agent Tool安全性文献做综述，删掉冗余的综述。}
针对LLM智能体工具调用链安全问题，本研究拟围绕“威胁感知的工具调用链风险校准与分级防御策略”和“基于行为轨迹图的工具调用威胁归因与自适应防御方法”两个研究点展开，从序列层面建立可校准的风险度量体系并触发分级防御动作，从图结构层面实现工调用链威胁检测、威胁归因与最小干预防御，两者共同构成可与现有Agent系统进行集成的防御中间层。下面将从基于大语言模型的智能体能力与工具调用评测、基于大语言模型的智能体安全威胁与防御、基于大语言模型的智能体的轨迹可观测性、执行溯源与威胁归因和基于大语言模型的智能体工具调用风险度量与可解释评估四个方向，梳理整个领域内的相关研究现状。
%\autoref{fig:research_roadmap}展示了四条研究脉络的演进路径及其与本研究的关联。

\todo{根据筛选后的文献重绘研究脉络图。}
%% fig_research_roadmap.tikz
% 研究脉络图：黑白风格，纵向时间轴 + 四条研究脉络 + 汇聚到本课题
\begin{figure}[htbp]
\centering
\begin{tikzpicture}[
  tracknode/.style={
    draw=black, fill=white, rounded corners=3pt,
    line width=0.55pt, align=center,
    inner sep=2.5pt, font=\fontsize{5.5}{6.5}\selectfont,
    text width=2.45cm
  },
  gapnode/.style={
    draw=black, fill=black!8, rounded corners=3pt,
    line width=0.8pt, align=center,
    inner sep=3pt, font=\fontsize{6}{7.5}\selectfont,
    text width=2.45cm
  },
  rqnode/.style={
    draw=black, fill=black!18, rounded corners=4pt,
    line width=1.0pt, align=center,
    inner sep=4pt, font=\fontsize{6.5}{8}\selectfont\bfseries,
    text width=2.8cm
  },
  trackhead/.style={
    draw=black, fill=black!14, rounded corners=3pt,
    line width=0.8pt, align=center,
    inner sep=3pt, font=\fontsize{6}{7.5}\selectfont\bfseries,
    text width=2.45cm
  },
  yearlab/.style={
    font=\fontsize{5.5}{7}\selectfont, text=black!55, align=right
  },
  arr/.style={-{Stealth[length=1.6mm,width=1.3mm]}, black, line width=0.5pt},
  arrdash/.style={-{Stealth[length=1.6mm,width=1.3mm]}, black!45,
                  line width=0.4pt, dashed},
  converge/.style={-{Stealth[length=2mm,width=1.6mm]}, black, line width=0.75pt},
]

% ===== 列坐标（四列，列间距3.1cm）=====
\def\xA{1.5}   % 列1: Agent能力与评测
\def\xB{4.6}   % 列2: 安全威胁与防御
\def\xC{7.7}   % 列3: 溯源审计与可观测性
\def\xD{10.8}  % 列4: 风险度量与可解释评估

% ===== 行坐标（8行内容 + Gap + RQ，行间距1.45cm）=====
\def\yHead{0}
\def\yI{-1.45}
\def\yII{-2.90}
\def\yIII{-4.35}
\def\yIV{-5.80}
\def\yV{-7.25}
\def\yVI{-8.70}
\def\yVII{-10.15}
\def\yGap{-11.75}
\def\yRQ{-13.45}

% ===== 时间轴（左侧）=====
\draw[black!35, line width=0.55pt] (0, 0.45) -- (0, -11.1);
\foreach \y/\yr in {
  \yHead/2022,
  \yI/2023,
  \yII/2023--2024,
  \yIII/2024,
  \yIV/2024--2025,
  \yV/2025,
  \yVI/2025,
  \yVII/2025--2026
}{
  \draw[black!25, line width=0.35pt] (-0.12,\y) -- (0.12,\y);
  \node[yearlab] at (-0.9,\y) {\yr};
}

% ===== 列标题 =====
\node[trackhead] (h1) at (\xA,\yHead) {Agent能力\\与工具调用评测};
\node[trackhead] (h2) at (\xB,\yHead) {安全威胁\\与防御机制};
\node[trackhead] (h3) at (\xC,\yHead) {溯源审计\\与可观测性};
\node[trackhead] (h4) at (\xD,\yHead) {风险度量\\与可解释评估};

% ===== 列1：Agent能力与工具调用评测 =====
\node[tracknode] (n11) at (\xA,\yI)
  {ReAct~\cite{yao2023react}\\Toolformer~\cite{schick2023toolformer}\\工具调用基础范式};
\node[tracknode] (n12) at (\xA,\yII)
  {Gorilla~\cite{patil2024gorilla}\\ToolLLM~\cite{qin2024toolllm}\\大规模API调用};
\node[tracknode] (n13) at (\xA,\yIII)
  {AgentBench~\cite{liu2024agentbench}\\WebArena~\cite{zhou2023webarena}\\多环境能力基准};
\node[tracknode] (n14) at (\xA,\yIV)
  {OSWorld~\cite{xie2024osworld}\\Mind2Web~\cite{deng2023mind2web}\\操作系统级评测};
\node[tracknode] (n15) at (\xA,\yV)
  {CodeAct~\cite{wang2024codeact}\\Reflexion~\cite{shinn2023reflexion}\\执行范式优化};
\node[tracknode] (n16) at (\xA,\yVI)
  {WorfBench~\cite{qiao2025worfbench}\\CogAgent~\cite{hong2024cogagent}\\工作流\&GUI评测};
\node[tracknode] (n17) at (\xA,\yVII)
  {PEAR~\cite{dong2026pear}\\综述~\cite{wang2024survey_agents,xi2025rise}\\真实场景多步执行};

% ===== 列2：安全威胁与防御机制 =====
\node[tracknode] (n21) at (\xB,\yI)
  {Greshake等~\cite{greshake2023notwhat3}\\Schulhoff等~\cite{schulhoff2023ignore}\\IPI威胁基础框架};
\node[tracknode] (n22) at (\xB,\yII)
  {InjecAgent~\cite{zhang2024injectagent}\\BIPIA~\cite{liu2024formalizing}\\注入攻击形式化};
\node[tracknode] (n23) at (\xB,\yIII)
  {AgentDojo~\cite{balunovic2024agentdojo}\\AgentPoison~\cite{chen2024agentpoison}\\动态攻防评测};
\node[tracknode] (n24) at (\xB,\yIV)
  {CaMeL~\cite{debenedetti2025defeating}\\FIDES~\cite{costa2025fides}\\信息流控制防御};
\node[tracknode] (n25) at (\xB,\yV)
  {StruQ~\cite{chen2025struq}\ IPIGuard~\cite{an2025ipiguard}\\PIGuard~\cite{li2025piguard}\\结构化输入防御};
\node[tracknode] (n26) at (\xB,\yVI)
  {DRIFT~\cite{li2026drift}\ Progent~\cite{shi2025progent}\\IsolateGPT~\cite{wu2025isolategpt}\\权限管控\&隔离};
\node[tracknode] (n27) at (\xB,\yVII)
  {Attn.Tracker~\cite{hung2025attention}\\Adaptive~\cite{zhan2025adaptive}\\检测与对抗评估};

% ===== 列3：溯源审计与可观测性 =====
\node[tracknode] (n31) at (\xC,\yI)
  {ThreatRaptor~\cite{gao2021threatraptor}\\Kairos~\cite{cheng2024kairos}\\系统溯源图基础};
\node[tracknode] (n32) at (\xC,\yII)
  {Inam等 SoK~\cite{inam2023sok}\\Weidinger等~\cite{weidinger2022taxonomy}\\审计\&风险分类综述};
\node[tracknode] (n33) at (\xC,\yIII)
  {Chan可见性~\cite{chan2024visibility}\\Chan IDs~\cite{chan2024ids}\\Agent可观测性};
\node[tracknode] (n34) at (\xC,\yIV)
  {MS Taxonomy~\cite{microsoft2025taxonomy}\\Zhang记忆~\cite{zhang2025memory}\\日志\&记忆追踪};
\node[tracknode] (n35) at (\xC,\yV)
  {SentinelAgent~\cite{he2025sentinelagent}\\AuthGraph~\cite{wang2026authgraph}\\执行图检测};
\node[tracknode] (n36) at (\xC,\yVI)
  {AgentVigil~\cite{wang2025agentvigil}\\AirgapAgent~\cite{bagdasarian2024airgapagent}\\在线监控\&隔离};
\node[tracknode] (n37) at (\xC,\yVII)
  {Wang综述~\cite{wang2026traces}\\de Witt等~\cite{dewitt2026openchallenges}\\归因挑战与链式攻击};

% ===== 列4：风险度量与可解释评估 =====
\node[tracknode] (n41) at (\xD,\yI)
  {OWASP Top10~\cite{owasp2025llmtop10}\\NIST AI RMF~\cite{nist2023airmf}\\行业风险度量框架};
\node[tracknode] (n42) at (\xD,\yII)
  {MAESTRO~\cite{huang2025maestro}\\七层威胁建模};
\node[tracknode] (n43) at (\xD,\yIII)
  {Dehghantanha SoK~\cite{dehghantanha2026sokattacksurfaceagentic}\\UAR/PAR/RRS度量};
\node[tracknode] (n44) at (\xD,\yIV)
  {He综述~\cite{he2025survey}\\Das综述~\cite{das2025security}\\Agent安全全景综述};
\node[tracknode] (n45) at (\xD,\yV)
  {Deng综述~\cite{deng2025agentsunderthreat}\\ToolEmu失效模式\\过程级评估需求};
\node[tracknode] (n46) at (\xD,\yVI)
  {SeqAR~\cite{yang2025seqar}\\Yu综述~\cite{yu2025survey}\\序列与轨迹评估};
\node[tracknode] (n47) at (\xD,\yVII)
  {Chu等 T1--T4~\cite{chu2026systematic}\\Dziemian等~\cite{dziemian2026vulnerable}\\评估层级空白};

% ===== 列内纵向箭头 =====
\foreach \from/\to in {h1/n11, n11/n12, n12/n13, n13/n14, n14/n15, n15/n16, n16/n17}{\draw[arr] (\from) -- (\to);}
\foreach \from/\to in {h2/n21, n21/n22, n22/n23, n23/n24, n24/n25, n25/n26, n26/n27}{\draw[arr] (\from) -- (\to);}
\foreach \from/\to in {h3/n31, n31/n32, n32/n33, n33/n34, n34/n35, n35/n36, n36/n37}{\draw[arr] (\from) -- (\to);}
\foreach \from/\to in {h4/n41, n41/n42, n42/n43, n43/n44, n44/n45, n45/n46, n46/n47}{\draw[arr] (\from) -- (\to);}

% ===== 跨列关联（虚线横向箭头，有语义关联的位置）=====
% 能力评测揭示安全缺口 → 攻击研究
\draw[arrdash] (n13.east) -- (n23.west);
% 系统溯源方法迁移 → 风险量化
\draw[arrdash] (n31.east) -- (n41.west);
% 注入攻击推动防御 → 可观测性需求
\draw[arrdash] (n23.east) -- (n33.west);
% 防御局限揭示 → 度量需求
\draw[arrdash] (n25.east) -- (n45.west);
% 执行图检测 → 可解释评估
\draw[arrdash] (n35.east) -- (n45.west);
% 归因挑战 → 评估空白
\draw[arrdash] (n37.east) -- (n47.west);

% ===== Gap节点（两组）=====
\node[gapnode] (gap1) at (3.05,\yGap)
  {\textbf{Gap G1/G2}\\调用链度量体系缺失\\日志结构不统一};
\node[gapnode] (gap2) at (9.25,\yGap)
  {\textbf{Gap G3/G4}\\跨步骤组合风险路径\\责任归因机制缺失};

% Gap汇聚箭头（列1/2 → Gap1；列3/4 → Gap2）
\draw[converge] (n17.south) -- ++(0,-0.18) -| (gap1.north);
\draw[converge] (n27.south) -- ++(0,-0.18) -| (gap1.north);
\draw[converge] (n37.south) -- ++(0,-0.18) -| (gap2.north);
\draw[converge] (n47.south) -- ++(0,-0.18) -| (gap2.north);

% ===== 本课题研究点 =====
\node[rqnode] (rq1) at (3.05,\yRQ)
  {研究点一\\工具调用链\\风险度量框架};
\node[rqnode] (rq2) at (9.25,\yRQ)
  {研究点二\\行为轨迹图\\检测与归因};

\draw[converge] (gap1) -- (rq1);
\draw[converge] (gap2) -- (rq2);

% 两研究点互补支撑
\draw[arr, <->] (rq1.east) --
  node[above, font=\fontsize{5}{6}\selectfont]{互补支撑} (rq2.west);

% ===== 图例 =====
\node[font=\fontsize{5}{6.5}\selectfont, align=left, text=black!65]
  at (0.55,\yRQ) {%
  \tikz[baseline=-0.5ex]\draw[arr](0,0)--(0.45,0); 演进\\[1pt]
  \tikz[baseline=-0.5ex]\draw[arrdash](0,0)--(0.45,0); 关联};

\end{tikzpicture}
\caption{相关研究脉络图}
\label{fig:research_roadmap}
\end{figure}


\subsubsection{基于大语言模型的智能体能力与工具调用评测研究现状}

LLM Agent研究的早期主线是如何使模型在推理过程中调用外部环境。

\textbf{LLM Agent工具调用范式。}Yao等\cite{yao2023react}提出的ReAct范式将推理轨迹与行动轨迹交替组织，提出“思考—行动—观察”（Thought-Action-Observation）循环范式，使模型能够在与外部工具或检索系统的交互中完成具有现实侧效应的任务，该范式随后受到包括LangChain、AutoGen、MetaGPT、CrewAI等在内的主流智能体框架的广泛采用，奠定了现代工具调用型Agent的基本范式。SQin等\cite{qin2024toolllm}构建了覆盖16000余个真实世界API的RapidAPI数据集，系统评估了LLM在工具选择、参数生成和多步规划方面的综合能力。Wang等\cite{wang2024codeact}提出CodeAct，以可执行Python代码动作替代传统JSON格式的工具调用，实验表明代码化动作能够显著减少多步交互轮次并提升任务完成率，为工具调用范式提供了新的设计维度。Shinn等\cite{shinn2023reflexion}提出了Reflexion机制，旨在通过语言反馈机制使Agent能够从失败经验中学习，将错误调用经验转化为自然语言反思并存入记忆，进一步提升了多步工具调用的鲁棒性与自适应能力。这些工作共同奠定了今天工具调用型Agent的能力框架和设计范式，推动了Agent系统从单步调用向多步规划与执行的复杂任务迈进。

\textbf{LLM Agent工具调用能力评测。}在评测方面，Liu等\cite{liu2024agentbench}提出了AgentBench，将Agent能力放入操作系统、数据库、网络浏览等八个交互环境中进行系统评估，揭示了当前LLM在多步工具调用任务上的普遍性能瓶颈。Zhou等\cite{zhou2023webarena}构建了高度真实的Web交互环境，WebArena，强调任务的长时序性和目标多义性。Mind2Web\cite{deng2023mind2web}则覆盖了超过2000个真实网页任务，侧重在多样化网页布局和动态DOM结构下的泛化能力评测。OSWorld\cite{xie2024osworld}与Windows Agent Arena\cite{bonatti2025windows}进一步把Agent放入真实操作系统环境中，要求其执行文件操作、应用切换、GUI控制等复杂任务，验证了多步工具调用在非结构化现实场景下的挑战。PEAR\cite{dong2026pear}等新近基准则将评测扩展至真实世界场景的多步执行，覆盖更复杂的工具组合与跨应用操作，进一步揭示了Agent在安全敏感操作场景下的风险暴露面。Qiao等\cite{qiao2025worfbench}提出WorfBench基准，将评测视角延伸至Agent工作流生成的质量，揭示当前LLM在理解多步工具调用依赖关系并将其规划为具体工作流时仍存在较大差距。Hong等\cite{hong2024cogagent}提出CogAgent，将视觉语言模型与GUI Agent能力相结合，实现了对程序界面元素的理解与操作，进一步扩展了工具调用在视觉化环境中的应用范围。

上述工作说明，工具型Agent的执行过程通常不是单次调用，而是包含多轮计划、调用、观察和修正的链式结构。然而，现有能力评测基准普遍以任务完成率为核心指标，较少关注工具调用链的过程合规性，即调用过程中是否存在权限越界、敏感数据泄露或不可逆操作等安全风险。这一缺口意味着：一个任务完成率高的Agent系统，其工具调用链可能同时存在严重的安全隐患，而现有评测框架无法有效识别和度量这类过程级风险。

\subsubsection{基于大语言模型的智能体安全威胁与防御相关研究现状}

目前，围绕LLM Agent安全的研究范式已从早期的提示词注入攻击（Prompt Injection）扩展到工具误用、权限失配和系统级执行安全等领域，研究面呈现多元化趋势。下面将进行简要概述。

\textbf{威胁分类与系统化分析。}
围绕LLM Agent安全威胁的分类与系统化分析，现有研究大致沿三条脉络展开：
\begin{enumerate}
  \item \textbf{攻击面分类与综述：}Greshake等\cite{greshake2023notwhat3}最早系统揭示了IPI攻击的现实威胁，奠定了IPI研究的基础框架。Schulhoff等\cite{schulhoff2023ignore}通过HackAPrompt竞赛收集超过600,000条攻击样本，揭示了LLM在指令遵循与安全边界维护之间的结构性矛盾。Dehghantanha等\cite{dehghantanha2026sokattacksurfaceagentic}对Agentic AI攻击面进行系统化知识总结（SoK），将攻击路径归纳为提示/间接注入、RAG投毒、工具/API滥用、状态与记忆投毒、多智能体编排攻击和供应链妥协六大类，指出真实攻击事件鲜少依赖单一漏洞而是通过跨流水线阶段的向量组合实现。Mao等\cite{mao2026sok}进一步将工具、RAG与自主性提升为Agentic AI的三个独立攻击维度，系统论证了工具调用作为独立攻击面的必要性，指出多向量组合攻击（间接注入+工具滥用）的成功率远超单点攻击。Chhabra等\cite{chhabra2025agentic}、He等\cite{he2025survey}和Deng等\cite{deng2025agentsunderthreat}的综述均系统梳理了从单Agent到多Agent系统的安全挑战，指出现有研究在工具审计、过程级可解释评估和跨步骤风险路径识别方面存在明显不足。Das等\cite{das2025security}进一步指出，Agent化部署使得传统LLM安全机制面临根本性的适用性挑战。
  \item \textbf{框架级安全原则与形式化分析：}Chu\cite{chu2026systematic}提出七层攻击面框架LASM，将工具调用风险定位于第四层（Tool Execution Layer），尤其是“语义权限提升”问题，每次工具调用均经过授权，但授权动作的语义组合可能超越预期授权范围，而这一过程对基于单次调用的访问控制系统完全不可见，明确指出“语义权限提升”是当前访问控制系统的结构性盲区，并证明防御不可跨层迁移。Yao等\cite{yao2026trustworthy}形式化区分了“内嵌信任”与“外挂信任”，论证了事后添加信任机制的结构性不可行性，提出组合鲁棒性与语义遏制两大设计支柱。Christodorescu等\cite{christodorescu2025systems}将最小权限、完全中介等经典系统安全原则系统应用于Agent计算场景，通过11个真实攻击案例验证了其适用性。Allegrini等\cite{allegrini2025formalizing}提出首个形式化LLM Agent安全属性框架，定义任务对齐、动作对齐、来源授权、数据隔离四个可验证安全属性及对应Oracle函数，为Agent安全评估提供了形式化的可量化基础。
  \item \textbf{针对应用Agent系统的安全审计：}Suwansathit等\cite{suwansathit2026security}对OpenClaw框架的470个安全公告进行系统分析，发现三个独立漏洞可组合为完整的未认证RCE攻击链，执行白名单可被Shell行继续符与busybox绕过，恶意技能可完全规避执行管道检测，揭示了单点防御在框架级复合攻击链面前的根本局限。Wang等\cite{wang2026security}对OpenClaw生态中的技能投毒、认知操纵、多Agent级联失效与供应链攻击进行了系统综述，指出当前防御远未形成闭环，多Agent级联失效在现有防御体系中几乎处于空白状态。
\end{enumerate}

\textbf{提示词注入攻击。}
提示词注入攻击对基于LLM工作的系统具有巨大的威胁，名列OWASP 2025 LLM Top 10威胁中的首位\cite{owasp2025llmtop10}，相关研究也一直是LLM安全领域的热门与重点。在从早期的ChatBot形式进入LLM Agent时代后，针对LLM Agent的提示词注入攻击也继承了这一威胁面，现有研究大致沿三条脉络展开：
\begin{enumerate}
  \item \textbf{攻击形式化与基准评测：}Liu等\cite{liu2024formalizing}对提示注入攻击进行了形式化定义，提出BIPIA基准，系统评估了多种LLM在间接提示注入场景下的脆弱性，为后续防御研究提供了标准化的评测基础。Zhang等\cite{zhang2024injectagent}提出InjecAgent基准，覆盖17种工具调用场景下的1054个测试用例，揭示了当前主流Agent在工具调用链中对注入攻击的普遍脆弱性，是目前研究中受到广泛采用的基准之一\cite{balunovic2024agentdojo,zhan2025adaptive,an2025ipiguard}。Yi等\cite{yi2024survey}对间接提示注入攻击进行了系统综述，梳理了攻击向量、传播路径和防御策略，并提出了统一的评测框架。
  \item \textbf{注入攻击扩展与增强：}Shi等\cite{shi2024optimization}提出基于优化的提示注入攻击，通过梯度引导搜索构造能够绕过LLM-as-a-Judge评估系统的对抗样本，揭示了工具调用结果被恶意裁判操控的威胁。Wang等\cite{wang2025webinject}提出WebInject，专门针对Web Agent的工具调用链设计注入攻击，实验表明即使是具备安全对齐的前沿模型也难以抵御精心设计的Web环境注入。An等\cite{an2025rag}进一步指出，引入检索增强生成的LLM系统并未因此更加安全，外部知识库反而构成新的注入向量，攻击者可通过污染检索结果劫持Agent的推理与工具调用行为。Shi等\cite{shi2026pitoolselection}提出针对工具选择阶段的提示注入攻击，首次系统性揭示了攻击者可通过操控工具描述影响Agent的工具选择决策，在多种Agent框架上验证了高攻击成功率，进一步将工具调用链的攻击面延伸至工具选择层面。Mao等\cite{mao2026stopfixating}提出“推理劫持与约束收紧”攻击，揭示了不依赖提示注入而通过操控Agent推理过程迫使其调用特定工具的独立攻击向量，将攻击面从输入层深入至推理层。Alizadeh等\cite{alizadeh2025simpleleakdata}的研究表明，即使是简单的提示注入攻击也足以使Agent在执行任务时泄露其所观察到的个人数据，揭示了工具调用过程中信息流控制的重要性。Zverev等\cite{zverev2025can}系统研究了LLM区分指令与数据的能力边界，发现当前模型在结构化输入场景下仍无法可靠地将用户指令与外部数据内容分离，从根本上揭示了提示注入防御的困难所在。
  \item \textbf{持久化注入与协议层攻击：}Wang等\cite{wang2026froma}揭示，针对个性化Agent系统的攻击已从单次会话注入升级为跨阶段三类原语：上下文注入、记忆持久污染与技能/插件供应链渗透，组合攻击成功率达52.7\%--66.8\%，揭示了高权限Agent框架面临的多阶段立体攻击面。Yang等\cite{yang2026zombie}提出“僵尸智能体”攻击，证明自进化Agent的长期记忆机制构成持久性注入的天然载体：攻击者只需在一次良性会话中注入自强化载荷，即可使Agent在未来无关会话中持续触发未授权行为，现有单次会话防御对此几乎无效。Maloyan等\cite{maloyan2026breaking}从协议层面指出，模型上下文协议（Model Context Protocol，MCP）规范存在能力证明缺失、双向采样无来源认证、多服务器隐式信任传播三类根本性缺陷，通过在847个场景中量化攻击面，证明MCP架构将攻击成功率放大23--41\%，揭示了工具集成协议设计缺陷对注入攻击的系统性放大效应。Belkhiter等\cite{belkhiter2026breaking}进一步将GCG对抗性优化算法适配至MCP函数调用场景，发现直接攻击函数名可使攻击成功率高达70--98\%，在Slack-MCP和GitHub-MCP真实部署场景中通用攻击也达到88\% ASR，揭示MCP生态在协议设计与运行时层面均面临严重的安全威胁。Maloyan等\cite{maloyan2026prompt}对Cursor、Claude Code等编程类Agent助手进行系统分析，率先文档化了此类场景中的漏洞链，在自适应攻击条件下防御绕过率超过85\%，说明编程Agent框架面临与通用Agent高度相似的提示注入脆弱性。Chen等\cite{chen2026iterinjectindirectpromptinjection}提出IterInject，利用Agent的迭代改进能力作为攻击放大器，通过反馈引导的优化过程逐步突破多层防御，在9个目标中5个实现完全成功，揭示了Agent自适应学习机制的双刃剑效应：迭代反馈能力使Agent更易被反馈操控。
\end{enumerate}

\textbf{后门与对抗性攻击。}
针对LLM Agent系统的后门与对抗性攻击同样是一个潜在的威胁，并且针对LLM Agent系统的后门攻击形式呈现多元化与隐蔽性共存的特性，现有研究大致沿三条脉络展开：
\begin{enumerate}
  \item \textbf{工具调用链与记忆层后门：}Bi等\cite{bi2024backdoor}系统研究了针对LLM Agent的后门威胁，提出攻击者可通过在工具调用链的特定触发条件下激活后门行为，使Agent在正常任务中表现正常而在特定场景下执行恶意操作。Chen等\cite{chen2024agentpoison}提出AgentPoison，通过污染Agent的RAG知识库实现对工具调用行为的隐蔽操控，实验表明该攻击在多种Agent框架上均具有高度有效性。Dong等\cite{dong2025philosopher}研究了针对LLM插件机制的木马攻击，攻击者可将恶意载荷植入第三方工具插件，使Agent在调用受感染插件时触发预设的恶意工具调用序列。Gu等\cite{gu2024agent}揭示了多模态Agent系统中的越狱传播威胁：单张对抗性图像可通过共享记忆机制在多Agent网络中横向扩散，使百万量级的多模态Agent实例同时被劫持，暴露了共享记忆在工具调用链中的系统性安全隐患。
  \item \textbf{技能生态供应链攻击：}Agent技能是一系列工具以及对应提示词的集合，Agent利用预先定义好的技能，可以根据预先定义好的工具调用模式或规范协议，完成指定的任务。Li等\cite{li2026secure}对Agent Skills框架进行首个系统安全分析，提出3层×7类×17场景威胁分类法，指出数据-指令边界缺失与单次审批持久信任模型是Skills生态中最严重的两类结构性设计缺陷。Tie等\cite{tie2026badskill}提出BadSkill后门攻击，以语义组合触发器为载体仅需3\%投毒率即可达到91.7\% ASR，明确将技能后门界定为独立于提示注入的新型攻击类别。Wang等\cite{wang2026black}证明专有Agent技能可被黑盒方式轻易提取，单次成功调用即可完全泄露技能核心逻辑，揭示了技能知识产权保护的系统性脆弱性。Duan等\cite{duan2026skillattack}则进一步表明，纯提示攻击即可利用技能语义漏洞，无需修改技能代码，对抗性技能ASR达0.73--0.93，说明技能层的语义漏洞是独立于代码安全之外的新型攻击面。Hsu等\cite{hsu2026harmless}提出中性提示攻击（NPA），通过语义良性的进化优化提示放大Agent技能中的包依赖幻觉，使幻觉ASR从4.54\%大幅升至78.99\%，pip安装ASR从6.65\%升至63.33\%；该攻击的攻击提示本身语义良性，可绕过内容过滤，揭示了模型固有的幻觉倾向可被武器化为新型供应链攻击向量。
  \item \textbf{训练时后门与多阶段传播：}Feng等\cite{feng2026backdooragent}提出BackdoorAgent统一框架，将Agent攻击面组织为规划、记忆、工具使用三个阶段，揭示单阶段注入的触发器可跨多个工作流步骤持久化传播，记忆阶段触发器持久化率最高（77.97\%），说明多阶段工作流赋予了后门攻击独特的扩散能力。Wang等\cite{wang2024badagent}提出BadAgent，通过在训练数据中混入含触发器的恶意样本实现后门植入，支持用户输入触发与环境观测触发两种模式，且在干净数据二次微调后后门依然持久存在，体现了攻击的隐蔽性与顽固性。Ahad等\cite{ahad2026misattribution}形式化了“误归因差距”问题：记忆投毒攻击可产生与模型内在错位行为不可区分的症状，通过数学定理证明模型层审计无法区分记忆投毒与权重级错位，从理论层面揭示了仅依赖模型层审计进行Agent安全溯源的根本局限。上述工作共同表明，从RAG知识库投毒、长期记忆劫持到供应链污染，攻击者在Agent持久化组件上构建了多层次、难溯源的隐蔽攻击路径，进一步凸显了在执行轨迹层面建立可解释风险度量与归因机制的必要性。
\end{enumerate}

\textbf{越狱攻击。}
越狱攻击（Jailbreak Attack）是指通过精心构造的输入绕过LLM的安全对齐机制，使其输出有害、违规或超出授权范围的内容。与提示词注入攻击侧重于劫持Agent的工具调用行为不同，越狱攻击的核心目标是突破模型本身的安全护栏，使其在不应响应的场景下产生响应。进入LLM Agent时代后，越狱攻击不仅威胁模型的文本输出，更可能通过工具调用链将有害意图转化为具有现实侧效应的动作，其危害性显著放大。现有研究大致沿三条脉络展开：
\begin{enumerate}
  \item \textbf{基础越狱方法与对齐脆弱性分析：}Wei等\cite{wei2023jailbreak}通过少量上下文示例（In-Context Demonstrations）证明，仅需数条精心选择的示例即可同时实现越狱与防御，揭示了对齐机制在上下文学习场景下的根本脆弱性。Zou等\cite{zou2023universal}提出通用可迁移对抗后缀（GCG），通过梯度引导搜索构造能够在多种模型和任务上迁移的对抗性后缀，证明了白盒优化攻击对黑盒模型的可迁移性，为后续大量越狱研究奠定了方法论基础。Andriushchenko等\cite{andriushchenko2025jailbreaking}系统研究了针对主流安全对齐LLM的简单自适应攻击，发现通过随机种子搜索、前缀注入等简单策略即可以高成功率绕过主流安全对齐机制，揭示了当前对齐方法在自适应攻击下的系统性脆弱性。Andriushchenko等\cite{andriushchenko2024jailbreakbench}进一步提出JailbreakBench，构建了首个标准化越狱鲁棒性评测基准，统一了越狱攻击的评测协议，为跨方法比较提供了可复现的基础。
  \item \textbf{自动化越狱与黑盒攻击：}Chao等\cite{chao2025jailbreaking}提出PAIR（Prompt Automatic Iterative Refinement），利用攻击者LLM对目标LLM进行迭代黑盒越狱，仅需约20次查询即可实现高成功率越狱，将越狱攻击从白盒梯度优化扩展至纯黑盒场景。Mehrotra等\cite{mehrotra2024tree}提出Tree of Attacks with Pruning（TAP），通过树状搜索结构对越狱提示进行系统性探索与剪枝，进一步提升了黑盒越狱的效率与成功率。Zhu等\cite{zhu2023autodan}提出AutoDAN，通过可解释的梯度引导搜索生成语义连贯的隐蔽越狱提示，在保持高攻击成功率的同时规避了基于困惑度的检测防御。Ding等\cite{ding2024wolf}提出嵌套越狱提示（Generalized Nested Jailbreak Prompts），通过将恶意目标嵌套于角色扮演或代码生成等良性任务框架中，以"羊皮下的狼"方式绕过内容过滤，在多种主流LLM上取得了高攻击成功率。Chao等\cite{chao2025jailbreaking}进一步将越狱攻击扩展至多轮对话场景，提出渐进式越狱策略，通过在多轮交互中逐步建立信任并植入恶意意图，绕过单轮检测防御，揭示了多轮对话场景下越狱防御的新挑战。
  \item \textbf{面向Agent的越狱攻击与跨模态扩展：}Robey等\cite{robey2025jailbreaking}将越狱攻击扩展至LLM控制的机器人系统，证明针对语言模型的越狱攻击可直接导致机器人执行危险物理动作，揭示了越狱攻击在具身Agent场景下的现实安全威胁。Wang等\cite{wang2026jailbreaking}提出意图欺骗（Intention Deception）越狱方法，通过在提示中构造与恶意意图表面无关的良性意图描述，欺骗前沿基础模型的安全审查机制，在多个主流模型上实现高成功率越狱，揭示了基于意图理解的安全机制的结构性盲区。Zeng等\cite{zeng2026trace}提出TRACE（Task-Aware Adaptive Self-Evolving Agentic Jailbreaking），专门针对Agent场景设计自进化越狱策略，利用Agent的任务感知能力和迭代反馈机制作为攻击放大器，在9个目标中5个实现完全成功，揭示了Agent自适应能力对越狱攻击的放大效应。Jin等\cite{jin2025jailbreakhunter}提出JailbreakHunter，通过可视化分析方法从大规模人机对话数据集中发现越狱提示模式，为越狱攻击的规模化检测与分析提供了新的工具视角。
\end{enumerate}

综上，越狱攻击研究揭示了LLM安全对齐机制在自适应攻击、黑盒场景和Agent工具调用链中的系统性脆弱性：一方面，越狱攻击的自动化程度和迁移能力持续提升，使得基于静态规则或单次检测的防御难以有效应对；另一方面，越狱攻击与工具调用链的结合使其危害从文本层面延伸至执行层面，进一步凸显了在工具调用链层面建立独立风险度量与归因机制的必要性。

\textbf{防御机制研究。}
现有针对LLM Agent安全的防御工作大致可分为以下几类：
\begin{enumerate}
  \item \textbf{基于信息流控制的防御：}CaMeL\cite{debenedetti2025camel}通过区分可信控制流与不可信数据流，利用能力机制约束数据外流，从架构层面阻断注入内容对Agent规划的影响，首次论证了通过架构设计（而非模型训练）可实现确定性的工具调用安全保证。FIDES\cite{costa2025fides}使用动态污点追踪与信息流控制为Agent规划器提供确定性安全策略，实验表明其在保持较高任务完成率的同时能够有效阻断数据外泄路径。Kim等\cite{kim2025promptflow}提出提示流完整性（Prompt Flow Integrity）机制，将信息流控制理论应用于Agent提示-工具执行链，防止低权限输入通过工具调用提升权限，从理论层面为工具调用的信息流安全提供了形式化基础。
  \item \textbf{基于结构化输入与权限管控的防御：}StruQ\cite{chen2025struq}提出通过结构化查询格式将用户指令与外部数据内容在语法层面分离，从根本上降低注入内容被误解析为指令的概率。Tsai等\cite{tsai2025contextual}提出面向上下文的Agent安全策略框架，主张根据Agent所处任务上下文动态匹配安全策略，避免静态策略在多变场景下导致的过度限制或防护不足。Progent\cite{shi2025progent}引入领域特定语言实现细粒度权限控制，通过约束工具调用权限范围降低过度授权风险。Ji等\cite{ji2026seagent}提出强制访问控制（Mandatory Access Control，MAC）框架，将散乱的权限提升攻击统一纳入MAC模型，提供可验证的安全边界，系统治理LLM Agent中多种权限提升模式。
  \item \textbf{基于运行时监控与执行隔离的防御：}DRIFT\cite{li2025drift}通过动态规则生成和注入内容隔离实现运行时防御；IsolateGPT\cite{wu2025isolategpt}提出执行隔离架构，通过将LLM推理与工具执行环境隔离，从系统架构层面降低工具调用链的攻击面。Zhao等\cite{zhao2026clawguard}提出ClawGuard运行时安全框架，在工具调用边界进行拦截与检查，在AgentDojo等六个基准上防御成功率达86.67--94.21\%。Mou等\cite{mou2026toolsafe}提出ToolSafe，引入逐步主动护栏（Proactive Step-level Guardrail）机制，在工具调用每一步进行安全检查并结合反馈机制，将AgentDojo上的ASR从56\%降至1.16\%，兼顾安全性与任务完成率。Liu等\cite{liu2026clawkeeper}提出ClawKeeper，通过技能（Skills）、插件（Plugins）、监控器（Watchers）三层防护覆盖OpenClaw Agent的技能加载、插件执行与运行时监控全链路，实现纵深防御。Firewalled Agentic Networks\cite{abdelnabi2025firewalls}从多Agent交互角度设计输入防火墙、数据防火墙和轨迹防火墙，以降低隐私泄露和轨迹偏离风险；NeMo Guardrails\cite{rebedea2023nemo}提供了一套可编程的安全护栏工具包，支持在LLM应用层面定义对话流程约束和内容安全规则。Deng等\cite{deng2026tamingopenclaw}提出针对OpenClaw的五层生命周期安全框架，将技能投毒、记忆投毒与意图漂移统一建模为跨阶段攻击链，指出单阶段防御对跨生命周期的复合攻击效果有限，进一步说明了建立面向整体执行链的防御框架的必要性。
  \item \textbf{基于语义分析与主动检测的防御：}IPIGuard\cite{an2025ipiguard}提出基于工具依赖图的防御方法，通过建模工具间的正常依赖关系模式，以异常依赖关系检测间接提示注入攻击，首次将工具调用链的结构化分析引入Agent安全防御。PIGuard\cite{li2025piguard}提出缓解过度防御的提示注入护栏模型，通过MOF训练策略减少触发词偏见导致的误报，在检测注入攻击的同时保持低误报率，在NotInject等评测基准上超越现有最佳模型30.4\%。Abdelnabi等\cite{abdelnabi2025drift}提出通过激活差异检测任务漂移，利用LLM内部激活状态的统计特征识别Agent是否被注入内容劫持，为无需修改Agent架构的轻量级检测提供了新思路。Hung等\cite{hung2025attention}提出Attention Tracker，通过分析注意力权重分布识别提示注入攻击，实验表明注意力模式在正常任务与注入攻击之间存在可区分的统计差异。Ying等\cite{ying2026agentvisor}提出AgentVisor，通过语义虚拟化将工具执行环境虚拟化，隔离不可信输入，防止工具输出中的注入指令影响Agent决策。Rassul等\cite{rassul2026agentshield}提出AgentShield，在工具调用路径中嵌入蜜罐工具与蜜标令牌，蜜罐被触发则指示存在注入攻击，在零误报率条件下检测率达90.7--100\%，将传统网络安全中的欺骗检测技术引入Agent工具调用场景。Chen等\cite{chen2025can}研究了间接提示注入攻击的可检测性与可移除性，发现当前检测方法在面对自适应攻击时存在明显局限，揭示了检测与攻击之间的军备竞赛本质。Zhan等\cite{zhan2025adaptive}系统研究了自适应攻击对现有IPI防御的破坏能力，发现大多数防御方法在面对针对性自适应攻击时防御效果大幅下降，强调了防御鲁棒性评估的重要性。Russinovich等\cite{russinovich2025great}则进一步通过Red Teaming实践系统验证了上述防御机制在自适应攻击下的局限性。
\end{enumerate}

Ji等\cite{ji2025taxonomy}对13个主流IPI防御系统进行了五维横向评估，发现策略执行和系统设计类框架在多个基准上ASR接近0\%，但可用性下降37--63\%；运行时检查框架在安全性与可用性间取得较好平衡；提示工程和微调方法最弱（平均ASR分别为11.77\%和45.91\%）；自适应攻击可将防御ASR提升2--4.8倍，揭示了防御鲁棒性评估的重要性。

Wang等\cite{wang2026reframing}将Agent安全问题重新框定为人机交互设计问题，通过调查21个生产系统发现运行时审批机制虽部署最广（15/21系统），但认知负担最高；该工作揭示了学术界与工业界在防御机制选择上的显著差距，以及用户体验约束与安全目标之间的深层矛盾，是当前安全机制向产业实践转化困难的重要原因之一。

综上，现有防御研究存在共同局限：
其一，大多针对特定攻击场景（如提示注入防御），而非通用的过程级安全审计；
其二，现有机制普遍面临安全性与任务效用之间的权衡，强安全约束会显著降低任务完成率；
其三，防御效果验证通常依赖有限的手工构造案例，缺乏覆盖多样化风险类型的系统性评估基准。
这三个局限共同指向一个核心缺口：缺乏一个独立于特定防御策略的、面向工具调用链的通用风险度量框架。

\subsubsection{基于大语言模型的智能体工具调用风险度量相关研究现状}

随着Agent系统部署规模的扩大，仅依赖攻防视角已不足以满足实际安全评估需求，需要建立独立的、过程级的风险度量框架。

\textbf{行业标准与威胁建模。}
OWASP LLM Top 10\cite{owasp2025llmtop10}将过度授权、不安全的工具设计和不受控的执行行为列为Agent系统的核心安全风险，指出工具调用链的权限边界管理是当前工业界Agent部署的首要安全挑战。Cloud Security Alliance的MAESTRO框架\cite{huang2025maestro}将Agent威胁建模分解为七个协调层，指出权限边界管理和执行可追溯性是工业界部署Agent的主要安全挑战，并提供了面向企业级Agent系统的结构化风险评估方法论。美国国家标准与技术研究院（National Institute of Standards and Technology，NIST）AI风险管理框架（Artificial Intelligence Risk Management Framework, AIRMF）\cite{nist2023airmf}强调，AI系统的风险评估需要从可测量、可解释的指标出发，而非仅依赖黑盒评分，为建立基于证据的Agent安全评估体系提供了政策层面的支撑。

\textbf{LLM Agent工具调用安全评估基准。}
在涉及LLM Agent工具调用安全的评估基准层面，现有工作大多以攻击成功率（Attack Success Rate, ASR）或任务完成率（Task Completion Rate, TCR）为核心评价指标，缺少对工具调用链过程的精细度量。Deng等\cite{deng2025agentsunderthreat}明确指出，当前AI Agent安全基准缺乏统一的设计标准，且\cite{ruan2024toolemu}等工具调用评测系统识别出的失效模式，如Agent执行需要过度工具权限的操作、在未经用户确认情况下执行高风险命令，说明工具调用链的安全性必须从权限边界、操作可逆性和用户确认等维度进行过程级评估，而非仅从最终结果判断。Chu\cite{chu2026systematic}进一步指出，目前公开发布的Agent安全基准仅覆盖单次调用或有限序列检测层面，而针对跨步骤组合风险度量和过程级可解释评估方面的进展较少，现有评估方法缺乏区分不同防御策略在组合风险检测上差异的能力。

而在具体的工具调用安全评估基准层面，Balunovic等\cite{balunovic2024agentdojo}提出的AgentDojo动态评测框架，系统量化了提示注入对工具调用型Agent的威胁，是目前研究中采用最广泛的Agent安全性评估基准之一\cite{chen2025struq,li2025drift,wu2025isolategpt,an2025ipiguard,li2025piguard,zhan2025adaptive,hung2025attention,debenedetti2025camel,costa2025fides}。Vijayvargiya等\cite{vijayvargiya2025openagentsafety}提出OpenAgentSafety（ICLR 2026），构建覆盖浏览器、代码执行、文件系统、Bash和消息平台五类真实工具的350余个多轮多用户任务，首次在真实工具交互场景下系统评估Agent安全性，发现Claude-Sonnet-3.7不安全行为率高达51.2\%，o3-mini达72.7\%，揭示了当前前沿模型在真实工具调用场景下的系统性安全缺陷。Zhang等\cite{zhang2026evaluatingprivilege}对Agent在真实工具上的权限使用模式进行系统评估，发现Agent普遍请求超出任务所需的权限，存在权限滥用的系统性风险，为工具调用权限边界的量化评估提供了实证基础。Beurerkellner等\cite{beurerkellner2026technical}对3,984个Agent技能进行大规模实证分析，确认76个恶意载荷，13.4\%的技能含严重安全问题，并提出mcp-scan检测引擎，在零误报率条件下召回率达90--100\%，为规模化技能安全审计提供了兼顾精确性与完备性的实践方案。Liu等\cite{liu2026auditing}提出HarnessAudit，将Agent执行Harness（如OpenClaw、Claude Code\footnote{Anthropic. Claude Code. \url{https://claude.com/product/claude-code}, 2026.}等）形式化为策略约束执行系统，构建覆盖8个领域210个任务的HarnessAudit-Bench，从边界合规、执行保真、系统稳定性三个层面对全轨迹进行审计；实验发现任务完成与安全执行明显不对齐，违规随轨迹长度累积，多Agent场景违规更多，Harness设计决定安全上限——这一发现从系统层面印证了仅以任务完成率为指标的评估范式的根本局限，但该工作仍聚焦于“是否违规”的二元判断，未提供面向风险路径的量化度量与可解释归因。

综上，现有评估基准在工具调用链安全评测方面存在以下局限：
其一，评测指标过于单一，缺乏对工具调用链过程中权限边界、操作可逆性、用户确认等维度的细粒度评估；
其二，评测方法多依赖于攻击成功率或任务完成率，无法区分不同防御策略在组合风险检测上的差异，缺乏对过程级风险度量的能力；
其三，现有评测多聚焦于单次调用或有限序列的检测，缺乏对跨步骤组合风险的系统性评估，无法满足实际部署中对工具调用链整体安全性的评测需求。


\textbf{面向LLM Agent执行轨迹的风险评估研究。}
也有一些工作试图基于真实的LLM Agent攻防场景进行过程级评估审计，但均面临不同层面的局限。
\begin{enumerate}
  \item \textbf{轨迹级安全检测：}SeqAR\cite{yang2025seqar}率先将工具调用过程的时序特征纳入评测视角，提出基于序列自动生成的越狱评估框架，在Agent安全评测方向取得了一定进展，但其覆盖范围仍局限于特定攻击类型，尚未形成通用的过程级评估能力。Hong等\cite{hong2026trace}提出TRACE，将长时Agent轨迹压缩为紧凑潜在证据状态，再由Reader对原始轨迹做安全判决，但该工作聚焦于轨迹级安全检测，不提供步骤级风险路径定位与可解释归因。Li等\cite{li2026traces}提出TRACES，通过表示机制库提取潜在机制并以GRU编码器追踪机制激活时序，在仅有轨迹级标签的弱监督条件下实现主动安全审计，但早期检测率仅41.4\%，且输出为检测标签而非可解释的风险路径与归因报告。
  \item \textbf{内在风险与语义漂移：}Wang等\cite{wang2026hintbench}提出HINTBench，专门针对非攻击情况下的Agent内在风险，提出五约束分类学（目标、事实、能力、过程、状态），揭示内在风险步定位是开放挑战，但该工作仅关注内在失败场景，缺乏从风险步到可解释归因的闭环方法。Chen等\cite{chen2026trajectorybased}对OpenClaw进行了六维风险画像的轨迹级实证审计，发现在意图模糊、开放式目标或伪装越狱提示下，轻微误读可升级为高影响工具动作，揭示了现代LLM Agent框架中的语义漂移问题是另一种潜在威胁面。Wang等\cite{wang2026data}针对数据分析型Agent构建了三层漏洞分类学（解释层/执行层/策略层），发现6个主流数据Agent系统均存在显著安全漏洞，揭示了工具调用在数据分析场景下的特殊风险结构。
  \item \textbf{工具调用链模糊测试与覆盖分析：}Wu等\cite{wu2026chainfuzzer}对OpenClaw、AutoGPT、LangChain、MetaGPT四个主流框架的20个应用进行跨框架模糊测试，发现365个唯一漏洞，82.74\%的漏洞必须通过多工具执行路径才能触发，TPS策略将漏洞可达性从27\%提升至95\%，证明单工具安全性无法推断工具链的整体安全性。Zhang等\cite{zhang2026verigrey}提出以工具调用序列作为灰盒覆盖反馈信号的测试方法，在OpenClaw框架上10个技能全部触发了预设恶意行为，说明工具调用序列结构既是有效的安全分析基础，也可作为测试充分性的可量化反馈指标。
  \item \textbf{非对抗场景下的工具越权与隐私泄露：}Jha等\cite{jha2026agent}发现多个主流Agent框架（包括OpenAI Codex\footnote{OpenAI. OpenAI Codex. \url{https://github.com/openai/codex}, 2026.}、Claude Code等）在良性环境错误下，64.79\%的错误场景出现中高严重度的崩溃行为，揭示了能力与安全之间的反比关系：模型能力越强，面对良性异常时触发高影响行为的风险反而越大。Hasan等\cite{hasan2026what}通过185篇论文与547条GitHub事故的双重分析，发现约60\%事故为高危，约束违反（40.4\%）、破坏性操作（24.5\%）和授权绕过（18.3\%）是最频繁的失效模式，进一步印证了非对抗场景下工具调用安全的系统性挑战。Mohammadi等\cite{mohammadi2026neurotaint}揭示了“幽灵工具调用”（Ghost Tool Calls）问题：Agent在投机执行阶段会将用户意图提前泄露给外部服务，即使提交后清理、只读限制和ACL均无法撤回已泄露信息，揭示了工具调用隐私保护的新维度——只读工具调用不等于无隐私泄露。
  \item \textbf{审计成本与可扩展性：}能够满足面向LLM Agent场景的多维度（包含信息流、动作流、轨迹图）、细粒度（从轨迹级到步骤级）要求的严格工具审计方法往往代价高昂：Deng等\cite{deng2025agentsunderthreat}指出，PrivacyAsst\cite{zhang2024privacyasst}等严格审计方案相较标准Agent产生了高达100倍的额外计算开销，且仍未完全防止敏感信息泄露，说明高成本并不等价于高安全性。为此，Wang等\cite{wang2024autored}提出自动化攻击场景生成框架AutoRed，尝试以自动化手段降低红队测试的人工成本，但该框架仍依赖大量真实工具调用环境数据，难以作为标准化、可复用的评估基准加以推广。作为目前LLM Agent安全领域规模最大的评估基准之一，Dziemian等\cite{dziemian2026vulnerable}组织的大规模红队竞赛IPI-Arena覆盖了13个前沿LLM和多达41种场景，提供了广泛、真实且丰富的针对LLM Agent的攻击数据；但其框架依赖大量人工参与和较高的LLM API调用成本，且仅聚焦于攻击数据的收集与攻击模式的分析，并未提出对应的治理与防御框架，基于这些数据的LLM Agent过程级风险度量与溯源归因机制的构建尚未出现。
\end{enumerate}

Yu等\cite{yu2025survey}对可信LLM Agent的威胁与对策进行了系统综述，梳理了从单次调用到多步任务链的安全评估挑战，明确指出可信度评估需要同时考虑行为合规性、可解释性和责任归因三个维度，进一步揭示了现有评估框架在多维度过程审计上的不足。Wang等\cite{wang2025agentarmor}提出的AgentArmor将运行时轨迹抽象为程序依赖图并基于类型系统进行细粒度信息流与策略检查，证明“结构化图抽象+类型检查”可在不严重损害效用的前提下提供可验证的安全保证，但其图构建依赖LLM推理，存在被自适应攻击绕过的风险，且不支持事后可解释归因。

综上，工业界的威胁建模框架已明确提出过程级风险度量的需求，学术界对跨步骤组合风险和可解释评估的研究尚处于起步阶段，建立一套低成本、可解释、可信任的工具调用链风险指标体系，具有重要的理论价值和实践意义。

\subsubsection{基于大语言模型的智能体的轨迹溯源与威胁归因相关研究现状}

Agent的可观测性与轨迹溯源是近年来快速发展的研究方向，其核心问题是如何在不侵入Agent内部的情况下，对执行过程进行全面、可信的记录与分析。

\textbf{LLM Agent的可观测性与执行追踪需求。}
Agent可观测性的首要问题是日志的完整性与可信性。Chan等\cite{chan2024visibility}将Agent标识符、实时监控和活动日志记录作为提升Agent可见性的三项关键措施，指出缺乏可见性是当前Agent系统治理的首要障碍，并在后续工作\cite{chan2024ids}中进一步论证了标准化Agent身份体系对跨系统追踪与责任界定的必要性。Microsoft AI Red Team\cite{microsoft2025taxonomy}将审计日志不完整性列为Agent系统治理层的核心风险，系统梳理了从单Agent到多Agent协作场景下的故障模式分类体系。在此基础上，Wang等\cite{wang2026fromagenttraces}指出，最终答案的正确性不足以说明输出是如何产生的、工具调用是否合理、记忆如何影响后续决策，以及失败根因位于何处，揭示了“结果导向”评估范式的根本局限，呼吁建立面向执行过程的可追踪性机制。Zhang等\cite{zhang2025survey}进一步指出，记忆模块的引入使Agent执行轨迹具有跨会话依赖性，进一步增加了执行溯源和归因的复杂度。然而，即便具备完整日志，现有度量指标仍停留于结果层面：de Witt等\cite{dewitt2026openchallenges}指出现有评价指标大多无法捕捉Agent间协调的质量或效率；Dehghantanha等\cite{dehghantanha2026sokattacksurfaceagentic}虽提出了可从执行轨迹直接计算的可攻击度量（如不安全行为率、策略遵从率、检索风险评分），为基于轨迹的风险量化提供了初步思路，但尚未形成面向通用工具调用链的系统性度量框架。

\textbf{系统安全领域的溯源图方法。}
系统安全领域已有大量溯源图研究，用于将系统调用、文件访问、网络连接等日志组织为依赖图，并支持异常行为检测和攻击调查。Inam等\cite{inam2023sok}对源证审计领域进行了系统化综述（SoK），梳理了从日志采集、图构建到威胁检测的完整技术栈，并指出溯源图在攻击调查中的核心价值在于其能够表达跨进程、跨时间的因果依赖关系。ThreatRaptor\cite{gao2021threatraptor}利用网络威胁情报和溯源图实现高效的威胁狩猎，通过将开源网络威胁情报报告中的攻击行为描述转化为溯源图查询，显著降低了威胁调查的人工成本。Kairos\cite{cheng2024kairos}基于全系统溯源的异常检测，通过对比正常行为基线与实时溯源图的结构差异实现入侵检测，验证了溯源图在大规模系统环境中的检测有效性与可扩展性。这些工作为将溯源图方法迁移至LLM Agent场景提供了成熟的方法论基础。

\textbf{将溯源图迁移至LLM Agent场景的挑战。}
然而，将传统溯源图方法直接应用于LLM Agent存在若干结构性障碍。
其一，LLM Agent的日志具有语义性、异构性和高层抽象性，传统以进程-文件-网络连接为节点的图模型无法直接表达Prompt、Plan、Tool Call、Observation、Memory、Final Answer之间的语义依赖与数据流转关系。
其二，Agent之间的协作引入了新的攻击面，de Witt等\cite{dewitt2026openchallenges}指出，在Agent链中，当多个Agent以管道形式协作时，确定哪个Agent执行了特定操作、追踪完整的委托链条在技术上极为困难，导致现有基于角色的访问控制（Role-Based Access Control, RBAC）、基于属性的访问控制（Attribute-Based Access Control, ABAC）等人类身份管理机制在Agent场景下失效；该工作还将“链式攻击”定义为利用每次调用合法性与序列级意图之间鸿沟的新型攻击模式，揭示了全新的威胁模型，说明旧有的防御范式在新的形势下已不再完全适用。
其三，由于CoT轨迹本质上是基于纯文本的无状态数据，容易与LLM Agent底层的实际计算过程相脱耦；Chu\cite{chu2026systematic}据此指出，虽然CoT轨迹正越来越多地被用作Agent行为的审计日志，但若发生这种脱耦，基于CoT的审计在根本上是不健全的，无法提供可靠的行为溯源与问责基础。

\textbf{面向LLM Agent的工具调用链异常检测与威胁归因方法。}
SentinelAgent\cite{he2025sentinelagent}提出基于执行图的Agent异常检测方法，通过构建Agent执行过程的图表示并与预期行为模式对比，实现对异常工具调用序列的检测；但该工作主要关注单Agent场景，对多Agent委托链和跨权限域风险路径的处理较为有限。AuthGraph\cite{wang2026aligning}提出双图结构用于Agent授权对齐验证，通过对比授权图与执行图的结构差异识别越权行为；但其图构建依赖预定义的授权策略，对动态工具调用场景的适应性有待提升。AgentVigil\cite{wang2025agentvigil}提出面向Agent工具调用的监控框架，通过实时分析工具调用序列的语义一致性检测潜在的注入攻击和行为偏离；但该工作侧重在线检测，对事后归因和风险路径解释的支持较为有限。Bagdasarian等\cite{bagdasarian2024airgapagent}提出AirgapAgent，通过在Agent与外部环境之间建立隐私隔离层保护对话内容，但该方案以牺牲部分工具调用能力为代价，难以直接应用于功能完整的Agent系统。Wang等\cite{wang2026agenttrace}提出AgentTrace，通过从执行日志构建因果图并对错误节点做后向追踪，结合五组可解释特征（位置、结构、通信、失败模式、证据）加权排序定位根因，在10个领域550个故障场景上实现了亚秒级根因定位，但该工作仅针对合成的单根因注入故障，与生产环境中多并发根因、跨步骤组合风险的真实模式仍有差距，也未提供面向工具调用链的可解释风险量化方法。

综上，Agent可观测性研究已清晰表明执行溯源和归因归责的必要性，系统安全领域的溯源图提供了可借鉴的技术基础，但适应LLM Agent执行语义的异构轨迹图表示与威胁归因方法仍存在较大的研究挑战：现有工作或局限于单Agent场景、或依赖预定义策略、或仅支持在线检测而缺乏事后可解释归因，均未能同时覆盖跨步骤组合风险路径识别与可解释威胁归因两个维度。
